\section{Related Works}
\textbf{Scalable Continuous Treatment Effect Estimation.}
Using neural networks to provide scalable solutions for estimating the effects of continuous-valued interventions has received significant attention in recent years. 
\cite{bica2020estimating} provide a Generative Adversarial Network (GAN) approach. 
The dose-response network (DRNet) \citep{schwab2020learning} provides a more direct adaptation of the TarNet \citep{shalit2017estimating} architecture for continuous treatments.
The varying coefficient network VCNet \citep{nie2021vcnet} generalizes the DRNet approach and provides a formal result for incorporating the target regularization technique presented by \cite{shi2019adapting}.
The RieszNet \citep{chernozhukov2021omitted} provides an alternative approach for targeted regularization.
Adaptation of each method is straightforward for use in our sensitivity analysis framework by replacing the outcome prediction head of the model with a suitable density estimator.

\textbf{Sensitivity and Uncertainty Analyses for Continuous Treatment Effects.}
The prior literature for continuous-valued treatments has focused largely on parametric methods assuming linear treatment/outcome, hidden-confounder/treatment, and hidden-confounder/outcome relationships \citep{carnegie2016assessing, dorie2016flexible, middleton2016bias, oster2019unobservable, cinelli2020making, cinelli2020omitted}. 
In addition to linearity, these parametric methods need to assume the structure and distribution of the unobserved confounding variable(s).
\cite{cinelli2019sensitivity} allows for sensitivity analysis for arbitrary structural causal models under the linearity assumption.
The MSM relaxes both the distributional and linearity assumptions, as does our CMSM extension. 
A two-parameter sensitivity model based on Riesz-Frechet representations of the target functionals, here the APO and CAPO, is proposed by \cite{chernozhukov2021omitted} as a way to incorporate confidence intervals and sensitivity bounds.
In contrast, we use the theoretical background of the marginal sensitivity model to derive a one-parameter sensitivity model.
\cite{detommaso2021causal} purport to quantify the bias induced by unobserved confounding in the effects of continuous-valued interventions, but they do not present a formal sensitivity analysis.
Simultaneously and independently of this work, \cite{marmarelisbounding} are deriving a sensitivity model that bounds the partial derivative of the log density ratio between complete and nominal propensity densities. 
Bounding the effects of continuous valued interventions has also been explored using instrumental variable models \citep{kilbertus2020aclass, hu2021generative, padh2022stochastic}.
